\section{Contexto e Relevância}

O mapeamento e a compreensão do Sistema Solar evoluíram de registros rudimentares do movimento dos astros para análises sofisticadas de sua dinâmica e história. Áreas como Astronomia Fundamental e Mecânica Celeste fornecem modelos matemáticos e simulações que revelam os caminhos que levaram à configuração planetária atual. Nas últimas décadas, técnicas numéricas como a integração direta das equações de movimento permitiram esclarecer processos antes apenas sugeridos teoricamente --- entre eles, a migração dos planetas gigantes e a reorganização das regiões de ressonância, elementos do Modelo de Nice \citep{Tsiganis2009}. Esses avanços reconstituem um passado dinâmico e caótico que moldou o Sistema Solar tal como o observamos hoje. Do ponto de vista metodológico, a ciência pode ser entendida como um processo de decisão baseado em dados: hipóteses são confrontadas com observações, e a interpretação confiável exige métodos estatísticos que quantifiquem incertezas e permitam distinguir sinal de ruído \citep{Wall2003}.

Dentro desse panorama, o Sistema Solar Exterior surge como laboratório natural. Nele destacam-se os \textbf{Centauros}, que orbitam entre Saturno e Netuno, e os \textbf{Objetos Transnetunianos} (TNOs), nas imediações e além da órbita de Netuno. Esses corpos têm tamanhos típicos da ordem de centenas de quilômetros e distâncias de 15 a 100~UA do Sol; a combinação de pequeno porte e grande distância resulta em baixo brilho aparente e dificulta medições precisas. Ainda assim, por preservarem características primitivas, carregam pistas valiosas sobre a formação e a evolução do Sistema Solar. Até recentemente, propriedades físicas fundamentais --- tamanho, forma, massa, atmosfera --- só podiam ser inferidas por métodos indiretos (fotometria, espectroscopia, telescópios espaciais). Um tipo de observação terrestre, de custo relativamente baixo e acessível a telescópios modestos, passou a desempenhar papel decisivo: \textbf{as ocultações estelares}.

Ocultações estelares ocorrem quando um corpo do Sistema Solar passa diante de uma estrela distante, bloqueando temporariamente sua luz \citep{Knieling2024}. A técnica tornou-se uma das mais precisas para caracterização física desses objetos, superada apenas por encontros diretos com sondas espaciais. A análise da \textbf{curva de luz} --- a variação temporal do fluxo estelar --- permite estimar a posição do objeto, inferir sua forma projetada e identificar estruturas periféricas (satélites, anéis) \citep{Cazeneuve2023}; sua versatilidade vai de asteroides próximos a TNOs extremamente distantes \citep{Fraser2024}. Ocultações são ainda empregadas há décadas no estudo de atmosferas planetárias \citep{Zhu2022}. O crescimento do volume de dados gerados por campanhas observacionais, porém, tornou a identificação manual de eventos inviável e impulsionou a adoção de métodos computacionais, em particular técnicas de \textbf{aprendizado de máquina} (\textit{Machine Learning}, ML). Neste contexto, ML refere-se a algoritmos que, a partir de dados de exemplo, ajustam internamente seus parâmetros para realizar uma tarefa (por exemplo, decidir se uma curva contém ou não uma ocultação), em vez de seguir regras fixas programadas à mão. No contexto das ocultações, o ML tem sido usado para detecção confiável de eventos, mitigação de ruídos e caracterização de sinais \citep{Cazeneuve2023}: redes neurais convolucionais (como o ODNet) analisam séries temporais ou imagens da estrela alvo e de estrelas de referência, distinguindo ocultações reais de flutuações atmosféricas ou artefatos instrumentais; em buscas por TNOs em campos densos, redes residuais fazem triagem dos candidatos gerados por técnicas de empilhamento e deslocamento (\textit{shift-and-stack}), reduzindo detecções espúrias \citep{Fraser2024}. Processos Gaussianos têm sido explorados para modelar curvas de luz e extrair parâmetros temporais com estimativa nativa de incertezas \citep{Knieling2024, Deisenroth2019}.

\section{Problema de Pesquisa}

O avanço instrumental e a multiplicação de redes de observação (incluindo ciência cidadã) expandiram fortemente o volume de dados de ocultações. Em muitas campanhas, cada participante registra milhares de quadros, embora a ocultação ocupe apenas uma pequena fração desse conjunto \citep{Cazeneuve2023}. O desafio central é identificar rapidamente, nesse volume massivo, quais curvas de luz exibem um evento positivo. A tarefa é complexa: ruídos observacionais e variabilidades atmosféricas (nuvens, \textit{seeing}) podem produzir quedas ou picos espúrios que mimetizam uma ocultação \citep{Cazeneuve2023, Knieling2024}; observações frequentemente operam no limite instrumental, exigindo definições cuidadosas de detecção e critérios estatísticos para não confundir excursões de ruído com sinal \citep{Wall2003}; em buscas por TNOs em regiões densas, sinais podem ser artefatos da subtração de estrelas fixas \citep{Fraser2024}. A detecção rápida, precisa e escalável tornou-se um gargalo no fluxo de trabalho científico. Métodos que dispensam ML --- por exemplo, janelas de suavização --- tendem a gerar muitos falsos positivos ou a falhar em estrelas fracas (magnitude $\geq 10$), em que o sinal é facilmente encoberto pelo ruído. Técnicas de ML emergem como alternativa: redes neurais e modelos de regressão moderna (como Processos Gaussianos) já são usados para filtrar candidatos, reconstruir dados e apoiar levantamentos \citep{Deisenroth2019, Fraser2024, CarleoMLPhysSci}, e se apresentam como caminho natural para os desafios de escalabilidade e precisão na era do \textit{big data} em astronomia \citep{Cazeneuve2023}.

\section{Objetivos da Dissertação}

O objetivo principal desta dissertação é desenvolver e validar uma abordagem automatizada, baseada em técnicas de ML, para detectar de forma robusta e eficiente eventos de ocultação estelar em dados fotométricos. Os objetivos específicos são: (a)~implementar um \textbf{fluxo de trabalho em etapas} (\textit{pipeline}): aquisição dos dados, pré-processamento e extração de \textbf{características numéricas} (\textit{features}) a partir de cada curva de luz --- por exemplo, amplitude do fluxo, mínimo da série suavizada, distância entre centróides de um agrupamento em dois níveis, etc.; (b)~\textbf{treinar} e comparar \textbf{classificadores} de ML (Regressão Logística, Random Forest, XGBoost e CatBoost), isto é, algoritmos que recebem um vetor de características e respondem ``positivo'' ou ``negativo'' (ou uma probabilidade de evento), ajustando seus parâmetros a partir de curvas cuja classificação verdadeira é conhecida; (c)~avaliar o desempenho com métricas de classificação binária (acurácia, precisão, revocação, F1-score, área sob a curva ROC) em dados reais e sintéticos, permitindo interpretar a capacidade de generalização e a importância de cada variável.

Neste trabalho optou-se por classificadores baseados em árvores de decisão e em vetores de características extraídos de forma explícita a partir das curvas, em vez de redes convolucionais que operam diretamente sobre as séries temporais brutas (como o ODNet). As razões são a \textbf{interpretabilidade}: é possível verificar quais grandezas (features) mais influenciam a decisão do modelo, o que é valioso para análise científica; a reprodutibilidade com volumes de dados moderados; e a facilidade de integração em fluxos já existentes. A metodologia explora a capacidade desses modelos em identificar padrões típicos de ocultação e distingui-los de ruídos instrumentais, atmosféricos ou intrínsecos, com perspectiva de incorporar também redes neurais em trabalhos futuros \citep{Cazeneuve2023}. A \textit{pipeline} desenvolvida atua como ferramenta de organização, priorização e triagem em grandes bases observacionais: (1)~indica quais curvas de luz merecem atenção prioritária e (2)~estima a probabilidade de evento em curvas ruidosas, visando superar limitações atuais de precisão e velocidade e, no futuro, permitir implementações em tempo real a bordo de telescópios.

A hipótese de pesquisa que norteia o trabalho é a de que é possível obter desempenho comparável ao de abordagens manuais ou baseadas em redes convolucionais usando um fluxo baseado em características clássicas (estatísticas do fluxo, derivadas, testes entre quartis, distância entre centróides do K-means) e em \textbf{ensembles} de árvores (combinações de muitos modelos de árvore de decisão), com a vantagem adicional de interpretabilidade e controle sobre cada etapa. O Capítulo 5 retoma os objetivos específicos e discute em que medida os resultados obtidos corroboram ou refinam essa hipótese.

A dissertação está organizada da seguinte forma. O \textbf{Capítulo 2} apresenta a fenomenologia das ocultações estelares e o processamento dos dados astronômicos que originam as curvas de luz. O \textbf{Capítulo 3} expõe os fundamentos dos modelos de ML utilizados --- em linguagem acessível a leitores com formação em física ou astronomia ---: classificação \textbf{supervisionada} (aprendizado a partir de exemplos com rótulo conhecido), agrupamento (K-means) e os conceitos de validação e generalização. O \textbf{Capítulo 4} descreve a metodologia completa da \textit{pipeline}: aquisição e armazenamento dos dados, pré-processamento, extração de características e treinamento e avaliação dos classificadores. O \textbf{Capítulo 5} apresenta os resultados dos experimentos e a discussão, retomando os objetivos específicos e as métricas de avaliação. O \textbf{Capítulo 6} apresenta as conclusões do trabalho, com contribuições e limitações de forma explícita.
